% ===============================================================
% GRADE 7 -- Real-world multiple-choice problems
% ===============================================================
\subsection{Real-World Problems (Grade 7)}
\probset{7}
{\small Read the situation, pick a letter within 20 seconds, then check the Answer Key.\par}
\begin{gradebody}

\textbf{Term 1}

\mcq{Lola leaves a basin of water outside. After two sunny days the water level has dropped, though no one used it. What happened to the water particles?}
{They were destroyed by heat}{They gained energy and escaped as vapor}{They shrank in size}{They turned into air}
{B}{Evaporation: particles gain energy and leave the surface; they are not destroyed.}{E}

\mcq{The smell of frying tuyo spreads through the house faster on a hot afternoon than on a cool night. Best explanation?}
{Smell particles grow larger in heat}{There are more particles at night}{Hot air is heavier}{Particles move faster at higher temperature}
{D}{Higher temperature gives particles more kinetic energy, so diffusion is faster.}{E}

\mcq{Naphthalene balls in the aparador get smaller over months without leaving any liquid. Which change of state is this?}
{Melting}{Sublimation}{Condensation}{Evaporation}
{B}{Solid goes directly to gas.}{E}

\mcq{A student heats crushed ice and reads the thermometer every minute. It stays at \SI{0}{\celsius} for several minutes while the ice melts. Why?}
{Heat is being used to overcome attractions between particles}{Heat has stopped entering the ice}{The thermometer is broken}{The ice is releasing heat}
{A}{During a change of state, absorbed energy breaks attractions instead of raising temperature.}{D}

\mcq{A student tests whether the amount of fertilizer affects the height of pechay. Which should be kept the same?}
{Amount of fertilizer}{Nothing; change everything}{Type of soil and amount of water}{Height of pechay}
{C}{These are controlled variables; fertilizer is independent and height is dependent.}{A}

\mcq{You need exactly \SI{50}{mL} of vinegar for an experiment. Which tool is most accurate?}
{Drinking glass}{Graduated cylinder}{Teaspoon}{Beaker}
{B}{Graduated cylinders have fine volume markings.}{E}

\mcq{Three classmates measure a table as \SI{1.52}{m}, \SI{1.52}{m}, and \SI{1.53}{m}. The true length is \SI{1.80}{m}. Their results are}
{precise but not accurate}{accurate and precise}{accurate but not precise}{neither}
{A}{The readings agree with each other (precise) but are far from the true value (not accurate).}{D}

\mcq{Sugar dissolves easily in hot salabat but some stays at the bottom of a cold glass of the same drink. Which factor explains this?}
{Stirring}{Colour of the drink}{Temperature}{Size of the glass}
{C}{Most solids are more soluble at higher temperature.}{E}

\mcq{A vendor dissolves \SI{50}{g} of salt in \SI{450}{g} of water to make brine for daing. What is the percent by mass of salt?}
{11\%}{10\%}{90\%}{50\%}
{B}{$50 \div (50 + 450) \times 100 = 10\%$; divide by the mass of the \emph{solution}.}{A}

\mcq{A soft drink kept in the ref stays fizzier than one left in a hot jeepney. Why?}
{Heat adds \ce{CO2}}{Gases dissolve better in cold liquids}{Cold drinks contain more sugar}{Cold removes air from the bottle}
{B}{Gas solubility decreases as temperature increases.}{A}

\mcq{Stirring powdered juice in water}
{makes the powder dissolve faster}{changes the solvent}{increases how much powder can dissolve}{always makes the solution saturated}
{A}{Stirring affects the rate of dissolving, not solubility.}{A}

\mcq{A cleaning solution turns red litmus paper blue. Which household item could it be?}
{Calamansi juice}{Soap solution}{Soft drink}{Vinegar}
{B}{Bases turn red litmus blue; soap is basic.}{E}

\mcq{Under the HPO, the onion cells look slightly blurry. What should you adjust?}
{Stage clips}{Coarse adjustment knob}{Fine adjustment knob}{Eyepiece position}
{C}{Use only the fine knob on high power to avoid hitting the slide.}{E}

\mcq{A cell appears at the left edge of the field of view. To move it to the centre, move the slide}
{to the right}{to the left}{upward}{downward}
{B}{The image is reversed, so you move the slide toward the side where the image is.}{D}

\mcq{You switch from HPO back to LPO while viewing cheek cells. What do you expect?}
{The cells disappear}{Bigger, darker image}{Smaller image, wider and brighter field}{Same image}
{C}{Lower magnification shows more area and admits more light.}{A}

\mcq{Comparing a kangkong leaf cell with a human cheek cell, which structure is found only in the leaf cell?}
{Chloroplast}{Mitochondrion}{Nucleus}{Cell membrane}
{A}{Only plant cells (green parts) have chloroplasts.}{E}

\mcq{Wilted pechay becomes firm again after watering, mainly because water fills the}
{chloroplast}{cell wall}{nucleus}{central vacuole}
{D}{The full central vacuole pushes against the cell wall (turgor).}{A}

\mcq{Heart muscle cells work non-stop and need a lot of energy. Which organelle would be most numerous in them?}
{Mitochondrion}{Chloroplast}{Vacuole}{Cell wall}
{A}{Mitochondria release energy through cellular respiration.}{A}

\textbf{Term 2}

\mcq{A farmer grows many new kamote plants from stem cuttings of one plant. The new plants are}
{genetically identical to the parent}{genetically varied}{from two parents}{produced by meiosis}
{A}{Vegetative propagation is asexual, so it makes clones.}{E}

\mcq{A farmer wants a pechay variety with new trait combinations that might resist a new pest. Which method is better?}
{Stem cuttings}{Seeds from cross-pollination}{Budding}{Cloning the best plant}
{B}{Sexual reproduction produces genetic variation.}{A}

\mcq{A human egg cell normally has how many chromosomes, and why?}
{46, for growth}{92, to have extra}{23, so that the zygote will have 46}{23, because of mitosis}
{C}{Meiosis halves the number; fertilization restores 46.}{A}

\mcq{All the tawilis in Taal Lake together make up a}
{community}{population}{ecosystem}{biosphere}
{B}{A population is one species in one area.}{E}

\mcq{A rice field has rice, frogs, snakes, insects, water, soil, and sunlight. If you consider only the living things, you are describing a}
{population}{community}{biome}{organism}
{B}{A community is all the populations, without the abiotic factors.}{A}

\mcq{In a rice field: rice $\rightarrow$ rat $\rightarrow$ snake $\rightarrow$ eagle. If the rice stores \SI{10000}{J}, about how much energy reaches the snake?}
{\SI{10000}{J}}{\SI{10}{J}}{\SI{100}{J}}{\SI{1000}{J}}
{C}{About 10\% passes at each step: $10000 \rightarrow 1000 \rightarrow 100$.}{A}

\mcq{Farmers kill most of the snakes in their rice fields. What is the most likely immediate effect?}
{The rat population increases}{No change}{Rice yield increases}{The eagle population increases}
{A}{Removing a predator lets its prey multiply, which can harm the rice.}{A}

\mcq{Why are there far fewer eagles than rats in an area?}
{Eagles eat plants}{Rats are producers}{Eagles reproduce faster}{Less energy is available at higher trophic levels}
{D}{Energy decreases at each level, so fewer top predators can be supported.}{A}

\mcq{A tricycle runs along a straight road at a steady \SI{30}{km/h}. The net force on it is}
{forward}{zero}{backward}{upward}
{B}{Constant velocity means balanced forces (net force zero).}{A}

\mcq{Two kids pull a toy: \SI{15}{N} east and \SI{10}{N} west. What happens?}
{\SI{25}{N} east}{\SI{5}{N} west}{No motion}{\SI{5}{N} east}
{D}{Opposite forces subtract; the net force is in the direction of the larger one.}{E}

\mcq{A jeepney goes \SI{4}{km} north, then returns \SI{4}{km} south to its terminal. Its displacement is}
{\SI{8}{km}}{\SI{8}{km} north}{\SI{0}{km}}{\SI{4}{km} north}
{C}{It ends where it started; the distance, however, is \SI{8}{km}.}{E}

\mcq{A bus travels \SI{120}{km} in \SI{2}{h}. Its average speed is}
{\SI{240}{km/h}}{\SI{30}{km/h}}{\SI{122}{km/h}}{\SI{60}{km/h}}
{D}{$v = d/t = 120/2$.}{E}

\mcq{A jeepney's distance--time graph is flat from minute 10 to 15. What happened?}
{It stopped to load passengers}{It moved fastest}{It was speeding up}{It moved backward}
{A}{No change in distance means it was at rest.}{A}

\mcq{A car's speedometer stays at \SI{40}{km/h} as it goes around a rotunda. Which is true?}
{Both are changing}{Speed and velocity are both constant}{Velocity is constant; speed is changing}{Speed is constant; velocity is changing}
{D}{Velocity includes direction, which keeps changing on a curve.}{D}

\textbf{Term 3}

\mcq{Why are kaldero bodies made of aluminium but their handles of plastic?}
{Aluminium conducts heat; plastic insulates}{Aluminium is an insulator}{For decoration}{Plastic conducts better}
{A}{Metals conduct heat to the food; insulators protect your hand.}{E}

\mcq{On a sunny afternoon at a beach in Batangas, the wind blows from the sea toward the land because}
{the land cools faster}{the sea heats faster}{sea air is warmer}{the land heats faster and the air above it rises}
{D}{Rising warm air over land pulls in cooler sea air (sea breeze, convection).}{A}

\mcq{Why do many people wear light-coloured clothes in the hot season?}
{They absorb more heat}{They trap cold air}{They reflect more radiation}{They conduct heat out of the body}
{C}{Light colours reflect radiant energy; dark colours absorb it.}{E}

\mcq{A thermos keeps sinigang hot for hours because it reduces heat transfer by}
{convection only}{conduction only}{radiation only}{conduction, convection, and radiation}
{D}{The vacuum layer blocks conduction and convection; the shiny walls reflect radiation.}{D}

\mcq{After one earthquake, Town X near the epicenter reports intensity VII and Town Y far away reports intensity III. Its magnitude}
{is different in each town}{cannot be measured}{is the same for both towns}{is higher in Town Y}
{C}{Magnitude is one value per quake; intensity varies with distance.}{A}

\mcq{After strong shaking near the coast, the sea suddenly pulls back and exposes the seafloor. What should you do?}
{Wait for a text alert first}{Take photos}{Move to high ground immediately}{Collect stranded fish}
{C}{The sea retreating is a natural tsunami warning sign; evacuate at once.}{E}

\mcq{A road cut shows rock layers where the hanging wall moved \emph{up} relative to the footwall. The rocks were}
{pulled apart}{not stressed}{sliding sideways}{pushed together}
{D}{Hanging wall up means a reverse fault, caused by compression.}{A}

\mcq{The school bell rings for an earthquake drill. What is the correct action inside a classroom?}
{Duck, Cover, and Hold under a sturdy desk}{Stand near the windows}{Use the stairs at once}{Run outside immediately}
{A}{Protect yourself from falling objects until the shaking stops.}{E}

\mcq{Earth is closest to the Sun in January, yet December--February is the cooler period in the Philippines. Why?}
{The Northern Hemisphere is tilted away, so sunlight is more slanted}{Earth is farthest from the Sun in December}{The Sun is weaker in December}{Clouds block all sunlight}
{A}{The tilt, not the distance, controls how intense the sunlight is.}{D}

\mcq{Heavy, long rains in western Luzon from July to September are usually brought by the}
{sea breeze}{land breeze}{Amihan}{Habagat}
{D}{The southwest monsoon carries moist air from the sea.}{A}

\end{gradebody}
\probstats
